\section*{Backup: Algorithms and Derivations}

\begin{frame}{Backup: Paper A Double Heston AES Algorithm}
  \scriptsize
  \begin{block}{Inputs and initialization}
    Use \(M,T,r,S_0\), variance parameters \((\kappa_i,\gamma_{\nu_i},\bar\nu_i,v_{0,i})\), and correlations \(\rho_{1,3},\rho_{2,4}\). Set \(\Delta t=T/M\), \(\nu_0^{(1)}=v_{0,1}\), \(\nu_0^{(2)}=v_{0,2}\), and \(X_0=\ln S_0\).
  \end{block}
  \begin{block}{For \(j=0,\ldots,M-1\)}
    \begin{enumerate}
      \item Sample \(\nu_{j+1}^{(1)}\) and \(\nu_{j+1}^{(2)}\) from their CIR noncentral chi-square distributions.
      \item Compute constants \(c_0,\ldots,c_6\) once for the step size and model parameters.
      \item Update \(X_{j+1}\) using the almost-exact log-price formula.
      \item Recover \(S_j=\exp(X_j)\).
    \end{enumerate}
  \end{block}
  \keymessage{This is the Paper A implementation logic: exact CIR sampling for both variance factors, followed by an almost-exact log-price update.}
\end{frame}

\begin{frame}{Backup: Paper A Double Heston Constants and Update}
  \begin{columns}[T,onlytextwidth]
    \begin{column}{0.49\textwidth}
      \begin{block}{CIR sampling constants}
        {\tiny
        \[
        \nu_{j+1}^{(i)}=\bar c^{(i)}\chi^2_{\delta^{(i)}}(\bar\kappa^{(i)}),\quad i=1,2,
        \]
        \[
        \bar c^{(i)}=\frac{\gamma_{\nu_i}^2}{4\kappa_i}(1-e^{-\kappa_i\Delta t}),
        \]
        \[
        \bar\kappa^{(i)}=\frac{4\kappa_i\nu_j^{(i)}e^{-\kappa_i\Delta t}}{\gamma_{\nu_i}^2(1-e^{-\kappa_i\Delta t})},
        \quad
        \delta^{(i)}=\frac{4\kappa_i\bar\nu_i}{\gamma_{\nu_i}^2}.
        \]}
      \end{block}
      \begin{block}{Log-price constants}
        {\tiny
        \[
        \begin{aligned}
        c_0&=\left(r-\frac{\rho_{1,3}\kappa_1\bar\nu_1}{\gamma_{\nu_1}}-\frac{\rho_{2,4}\kappa_2\bar\nu_2}{\gamma_{\nu_2}}\right)\Delta t,\\
        c_3&=\frac{\rho_{1,3}}{\gamma_{\nu_1}},\quad c_4=\frac{\rho_{2,4}}{\gamma_{\nu_2}},\\
        c_5&=(1-\rho_{1,3}^2)\Delta t,
        \quad c_6=(1-\rho_{2,4}^2)\Delta t .
        \end{aligned}
        \]}
      \end{block}
    \end{column}
    \begin{column}{0.49\textwidth}
      \begin{block}{Remaining constants}
        {\tiny
        \[
        \begin{aligned}
        c_1&=\left(\frac{\rho_{1,3}\kappa_1}{\gamma_{\nu_1}}-\frac12\right)\Delta t-\frac{\rho_{1,3}}{\gamma_{\nu_1}},\\
        c_2&=\left(\frac{\rho_{2,4}\kappa_2}{\gamma_{\nu_2}}-\frac12\right)\Delta t-\frac{\rho_{2,4}}{\gamma_{\nu_2}} .
        \end{aligned}
        \]}
      \end{block}
      \begin{block}{Update}
        {\tiny
        \[
        \begin{aligned}
        X_{j+1}&=X_j+c_0+c_1\nu_j^{(1)}+c_2\nu_j^{(2)}+c_3\nu_{j+1}^{(1)}+c_4\nu_{j+1}^{(2)}\\
        &\quad+\sqrt{c_5\nu_j^{(1)}}\,(W_{j+1}^{(1)}-W_j^{(1)})\\
        &\quad+\sqrt{c_6\nu_j^{(2)}}\,(W_{j+1}^{(2)}-W_j^{(2)}).
        \end{aligned}
        \]}
      \end{block}
    \end{column}
  \end{columns}
\end{frame}

\begin{frame}{Backup: Paper B AEMS Algorithm}
  \scriptsize
  \begin{block}{Inputs and initialization}
    Use \(N,M,T,r,S_0,\kappa_1,\kappa_2,\xi_1,\xi_2,\rho_{12},\rho_{13},\rho_{23},\theta,v_0,v'_0\). Set \(v=v_0\), \(v'=v'_0\), \(X_0=\ln(S_0)\), and \(\Delta t=T/M\).
  \end{block}
  \begin{block}{For \(i=0,\ldots,M-1\)}
    \begin{enumerate}
      \item Generate independent Wiener increments \(\Delta W^{(1)},\Delta W^{(2)},\Delta W^{(3)}\).
      \item Sample \(v'_{i+1}=c\chi^2_\delta(\bar\kappa)\), where
      \[
      \delta=\frac{4\kappa_2\theta}{\xi_2^2},\quad
      c=\frac{\xi_2^2(1-e^{-\kappa_2\Delta t})}{4\kappa_2},\quad
      \bar\kappa=\frac{4\kappa_2v'_ie^{-\kappa_2\Delta t}}{\xi_2^2(1-e^{-\kappa_2\Delta t})}.
      \]
      \item Update \(v_{i+1}\) by the Milstein-style step.
      \item Update \(X_{i+1}\) using the correlation-decomposition formula and return \(S=e^X\).
    \end{enumerate}
  \end{block}
\end{frame}

\begin{frame}{Backup: Paper B AEMS Update Formulas}
  \begin{columns}[T,onlytextwidth]
    \begin{column}{0.49\textwidth}
      \begin{block}{Variance update}
        {\tiny
        \[
        \begin{aligned}
        v_{i+1}=v_i&+\kappa_1(v'_i-v_i)\Delta t
        +\xi_1\sqrt{v_i}\Delta W_i^{(2)}\\
        &+\frac14\xi_1^2\left((\Delta W_i^{(2)})^2-\Delta t\right).
        \end{aligned}
        \]}
      \end{block}
      \begin{block}{Correlation constants}
        {\tiny
        \[
        q=\frac{\rho_{12}-\rho_{13}\rho_{23}}{\sqrt{1-\rho_{23}^2}},\quad
        \rho^*=\sqrt{1-\rho_{13}^2-q^2},
        \]
        \[
        \rho^{**}=q-\rho_{13}\frac{\sqrt{1-\rho_{23}^2}}{\rho_{23}}.
        \]}
      \end{block}
    \end{column}
    \begin{column}{0.49\textwidth}
      \begin{block}{Log-price update}
        {\tiny
        \[
        \begin{aligned}
        X_{i+1}=X_i&+\left(r-\frac12v_i\right)\Delta t
        +\rho^*\sqrt{v_i}\Delta W_i^{(1)}
        +\rho^{**}\sqrt{v_i}\Delta W_i^{(2)}\\
        &+\rho_{13}\left(\frac{v_{i+1}-v_i-\kappa_1(v'_i-v_i)\Delta t}{\xi_1\rho_{23}}\right).
        \end{aligned}
        \]}
      \end{block}
      \begin{block}{AEMS-SOR}
        AEMS-SOR adds second-order correction terms to improve weak convergence. Keep the details in backup unless a committee member asks for the exact second-order expression.
      \end{block}
    \end{column}
  \end{columns}
\end{frame}

\begin{frame}{Backup: Paper C Brownian Projection and Conditional Step}
  \begin{columns}[T,onlytextwidth]
    \begin{column}{0.48\textwidth}
      \begin{block}{Projection coefficients}
        {\scriptsize
        \[
        \beta_2=\frac{\rho_{12}-\rho_{13}\rho_{23}}{1-\rho_{23}^2},\qquad
        \beta_3=\frac{\rho_{13}-\rho_{12}\rho_{23}}{1-\rho_{23}^2},
        \]
        \[
        \sigma_\perp^2=\frac{1-\rho_{12}^2-\rho_{13}^2-\rho_{23}^2+2\rho_{12}\rho_{13}\rho_{23}}{1-\rho_{23}^2}.
        \]}
      \end{block}
      \begin{block}{Brownian decomposition}
        {\scriptsize
        \[
        \dd W_t^{(1)}=\beta_2\dd W_t^{(2)}+\beta_3\dd W_t^{(3)}+\sigma_\perp\dd B_t^\perp.
        \]}
      \end{block}
    \end{column}
    \begin{column}{0.48\textwidth}
      \begin{block}{Conditional Gaussian representation}
        {\scriptsize
        \[
        X_{t_{n+1}}=y+Z_n+I_n+\sqrt{B_n}\xi,
        \qquad \xi\sim N(0,1),
        \]
        \[
        Z_n=\int_{t_n}^{t_{n+1}}\sqrt{v_s}(\beta_2\dd W_s^{(2)}+\beta_3\dd W_s^{(3)}),
        \]
        \[
        I_n=\int_{t_n}^{t_{n+1}}(r-\tfrac12v_s)\dd s,
        \qquad
        B_n=\sigma_\perp^2\int_{t_n}^{t_{n+1}}v_s\dd s.
        \]}
      \end{block}
    \end{column}
  \end{columns}
\end{frame}

\begin{frame}{Backup: Paper C Hybrid LSMC-PDE Recursion I}
  \scriptsize
  \begin{block}{Training paths and backward propagation}
    \begin{enumerate}
      \item Simulate \(N\) training paths of the variance factors on an Euler grid containing the Bermudan exercise dates.
      \item Store exercise-date volatility states and accumulate \(Z_n^j\), \(I_n^j\), and \(B_n^j\) over each interval \([t_n,t_{n+1}]\).
      \item Set \(\widehat V_M^N(s_i,v,w)=h_M(s_i)\).
      \item For \(n=M-1,\ldots,1\), compute \(\widehat C_n^j(s_i)\) using FFT propagation.
      \item For each asset-grid point \(s_i\), compute
      \[
      A_n^N=\frac1N\sum_{j=1}^N\phi(V_n^j)\phi(V_n^j)^\top,
      \qquad
      b_n^N(s_i)=\frac1N\sum_{j=1}^N\phi(V_n^j)\widehat C_n^j(s_i).
      \]
    \end{enumerate}
  \end{block}
\end{frame}

\begin{frame}{Backup: Paper C Hybrid LSMC-PDE Recursion II}
  \scriptsize
  \begin{columns}[T,onlytextwidth]
    \begin{column}{0.49\textwidth}
      \begin{block}{Regression and continuation}
        \begin{enumerate}
          \item Set \(a_n^N(s_i)=[A_n^N]^{-1}_R b_n^N(s_i)\).
          \item Fit the continuation surface:
          \[
          \widehat C_n^N(s_i,v,w)=a_n^N(s_i)\cdot\phi(v,w).
          \]
          \item Apply the exercise decision:
          \[
          \widehat V_n^N=\max\{h_n,\widehat C_n^N\}.
          \]
        \end{enumerate}
      \end{block}
    \end{column}
    \begin{column}{0.49\textwidth}
      \begin{block}{Time-zero estimator}
        \begin{enumerate}
          \item Average first-step pre-surfaces:
          \[
          \widehat C_0^N(s_i)=\frac1N\sum_{j=1}^N\widehat C_0^j(s_i).
          \]
          \item Interpolate at \(S_0\).
          \item Return
          \[
          \widehat V_0^N(S_0)=\max\{h_0(S_0),\widehat C_0^N(S_0)\}.
          \]
        \end{enumerate}
      \end{block}
    \end{column}
  \end{columns}
\end{frame}
